\documentclass{article}
\usepackage{ctex}
\usepackage{geometry,ragged2e,setspace,graphicx,verbatim}
\begin{document}
    \newgeometry{left=2.5cm,bottom=0cm}
    \begin{titlepage}
    \begin{center}
        \setstretch{2}
        \includegraphics{./50mun.png}\\
        \begin{flushleft}
            \vspace*{3em}
            \normalfont{\zihao{1}\CJKfamily{hei}\bfseries 合肥市五十中第十届模拟联合国大会}\\
            \normalfont{\bfseries \Large\textrm{The 10th Model United Nations of Hefei No. 50 Middle School}}\\
            \vspace*{6em}
            \centering\normalfont{\zihao{0}\CJKfamily{hei}\bfseries 背景文件}
            \\
            \vspace*{3em}
            \centering\normalfont{\Huge\bfseries \textrm{Background Guide}}\\
            \vspace*{5em}
            \centering\normalfont{\zihao{0}\CJKfamily{hei}\bfseries 联合国安全理事会}
            \\
            \vspace*{3em}
            \centering\normalfont{\Huge\bfseries \textrm{United Nations Security Council }}\\
            \vspace*{3em}
            \centering\normalfont{\Huge\bfseries\CJKfamily{hei} 议题：第四次中东战争的调停与善后}\\
            \vspace*{3em}
            \centering\normalfont{\Huge\bfseries\CJKfamily{hei} 作者：模拟联合国社团学术部}
        \end{flushleft}
    \end{center}
    \end{titlepage}
    \clearpage
    \tableofcontents
    \clearpage
    \section{委员会致辞}
    \clearpage
    \section{议题介绍}
        本次会议的议题是第四次中东战争的调停与善后，包括调停和善后两个方面。历史上1973年10月22日早上，联合国通过了停火决议。本次会议背景设置在停火之前，代表必须先撰写决议草案停战。决议草案的内容包括强制停战撤军、划定联合国管理区域等。该部分预计在第一议程完成。

        第二部分是善后，可以包括多方面，如难民、能源问题等。具体可以参考临时议程书。

        下面会介绍第四次中东战争的历史与现状。我们将会从第一次世界大战中东格局初步形成，讲到第四次中东战争以色列转败为胜。
        \subsection{历史与现实}
            \subsubsection{第一次世界大战：强弩之末的帝国、现代中东的形成}
                \subsubsection{战前中东}
                一战前的中东版图和现在大不相同，以色列、约旦、叙利亚、沙特阿拉伯等大部分区域处在一个庞大帝国的统治之下，
\end{document}